% Shared appendix: proofs of the results of the analytical part.
\section{Proofs}\label{app:proofs}
Throughout, $\Omega\subset\R^2$ is open and connected, $\mathcal A$ is $\mathcal D'(\Omega)$ or $C^\infty(\Omega)$, and
$\ring=\C[\tau,\xi]$ acts on $\mathcal A$ by $\tau\mapsto\partial_t$, $\xi\mapsto\partial_x$. For $z=(\tau_z,\xi_z)\in\C^2$,
$\mathfrak m_z=(\tau-\tau_z,\xi-\xi_z)$ is its maximal ideal and $\psi_z(t,x)=e^{\tau_zt+\xi_zx}$.

\subsection{Algebraic toolkit}
\begin{proof}[Proof of \cref{lem:gcd}]
\emph{Coprime polynomials meet in finitely many points.} Let $f,g\ne0$ be coprime. By Gauss's lemma they are
coprime in the principal ideal domain $\C(\xi)[\tau]$, so $af+bg=1$ there; clearing denominators,
$a'f+b'g=r(\xi)\ne0$ with $a',b'\in\ring$. At a common zero $r(\xi_0)=0$, which leaves finitely many $\xi_0$;
exchanging the variables, finitely many $\tau_0$. Hence $\Var(f,g)$ is finite.

\emph{Splitting.} Put $J=\{f/g:\ f\in I\}$ with $g=\gcdop(I)$; it is an ideal and $I=gJ$, and no prime divides all
of $J$. If $J\ne\ring$, pick $0\ne f\in J$ with distinct prime factors $\pi_1,\dots,\pi_r$ and, for each $a$,
$g_a\in J$ with $\pi_a\nmid g_a$. The sets $\{c\in\C^r:\pi_a\mid\sum_bc_bg_b\}$ are proper linear subspaces, so some
$c$ avoids all of them; then $f$ and $\sum_bc_bg_b\in J$ are coprime and $\Var(J)\subseteq\Var(f,\sum_bc_bg_b)$ is
finite. By the Nullstellensatz a power of $\sqrt J$ lies in $J$, and $J=\ring$ when $\Var(J)=\emptyset$.
\end{proof}

\begin{proof}[Proof of \cref{lem:ep}]
($\Rightarrow$) $Q(\partial)(q\,\psi_z)=\big(Q(z+\partial)q\big)\psi_z$ for every polynomial $q$, so the derivatives of
$q\,\psi_z$ stay in the finite-dimensional space of polynomials of degree $\le\deg q$ times $\psi_z$. Moreover
$\prod_k\mathfrak m_{z_k}^{\deg q_k+1}\subseteq\Ann(w)$, so $\Var(\Ann w)\subseteq\{z_k\}$; conversely
$Q\in\Ann(w)$ forces $Q(z_k+\partial)q_k=0$ by the linear independence of distinct exponentials, whose top-degree part
is $Q(z_k)$ times that of $q_k$, so $Q(z_k)=0$.

($\Leftarrow$) Let $\ring w$ be finite-dimensional and $p,q$ the minimal polynomials of $\partial_t$ and
$\partial_x$ restricted to $\ring w$, so that $p(\partial_t)w=0$ and $q(\partial_x)w=0$. On an open rectangle $I'\times J'\subset\Omega$, induction
on $\deg p$ shows that every distribution with $p(\partial_t)v=0$ is $\sum_{j,l}c_{jl}(x)\,t^le^{\lambda_jt}$ with
$c_{jl}\in\mathcal D'(J')$ and $\lambda_j$ the roots of $p$ (the base case uses that a distribution with
$\partial_tV=0$ on a rectangle is $c(x)\otimes1$). Applying $q(\partial_x)$ and testing against functions dual to
$t^le^{\lambda_jt}$ gives $q(\partial_x)c_{jl}=0$, so each $c_{jl}$ is an exponential-polynomial of $x$. Thus $w$ is an
exponential-polynomial $E$ near every point. The set where $w$ coincides near $y$ with a fixed $E$ is open, and it is
closed because two exponential-polynomials that agree on an open set agree everywhere; by connectedness it is
$\Omega$.
\end{proof}

\subsection{Identification}
\begin{proof}[Proof of \cref{thm:ident}]
Let $I=\Ann(u)\ni P$ and write $I=G\,J$ as in \cref{lem:gcd}. Since $P\in I$, $G\mid P$, so
$G=c'\prod_iP_i^{f_i}$ with $0\le f_i\le e_i$. First, $G\sim P$ iff $I=(P)$: if $G\sim P$ then
$I=PJ\subseteq(P)\subseteq I$; if $I=(P)$, $P$ divides all of $I$ and lies in $I$.

(a), $\Leftarrow$ by contraposition. If $I\ne(P)$, some $f_i<e_i$, so $P/P_i=G\,H$ with $H\in\ring$. A power
$\mathfrak n^{\ell}$ of the ideal $\mathfrak n$ of the finite set $\Var(J)$ lies in $J$, and for $h\in\mathfrak n^{\ell}$,
$h\cdot P/P_i=G(Hh)\in GJ=I$, i.e.\ $h(\partial)v_i=0$. So $\Var(\Ann v_i)\subseteq\Var(J)$ is finite and
$v_i\in\EP$ by \cref{lem:ep}.

(a), $\Rightarrow$ by contraposition. If $v_i\in\EP$, $\Var(\Ann v_i)$ is finite while $\Var(P_i)$ is infinite, so
some $h\in\Ann(v_i)$ is not divisible by $P_i$. Then $Q=h\cdot P/P_i$ satisfies $Q(\partial)u=h(\partial)v_i=0$; if
$Q=PA$, cancelling $P/P_i$ gives $h=P_iA$, a contradiction. So $Q\in I\setminus(P)$.

(b) $\Var(I)=\Var(G)\cup\Var(J)$ with $\Var(J)$ finite. If $f_i\ge1$, $\Var(P_i)\subseteq\Var(I)$; if $f_i=0$,
$\Var(P_i)\cap\Var(I)$ lies in the finite set $\bigcup_{f_j\ge1}\Var(P_i,P_j)\cup\Var(J)$, so
$\Var(P_i)\not\subseteq\Var(I)$. Hence $\Var(I)=\Var(P)$ iff all $f_i\ge1$, and $f_i=0$ iff $v_i'\in\EP$ by the
argument of (a) with $P/P_i^{e_i}$ in place of $P/P_i$.

(c) If $e_i=1$ then $v_i=v_i'$. If $P$ is monic of degree $m$ in $\tau$, divide any $Q\in\ring$ by $P$ in
$\C[\xi][\tau]$: $Q=AP+\sum_{j<m}\varphi_j(\xi)\tau^j$ uniquely. Then $Q\in I$ iff $\sum_j\varphi_j(\partial_x)\partial_t^ju=0$, and
$Q\in(P)$ iff all $\varphi_j=0$.
\end{proof}

\begin{proof}[Proof of \cref{prop:unique}]
(a) If $I=(P)$, every annihilator is $SP$ and $\deg SP\ge\deg P$ with equality iff $S$ is a nonzero constant.
(b) If $I=(P)$ then $G\sim P$ and $J=\ring$, so $\sigp(u)=\emptyset$. Conversely, $\Var(J)=\emptyset$ gives $J=\ring$ by
the Nullstellensatz, hence $I=(G)=(P)$. (c) Let $f_i$ be the multiplicity of $P_i$ in $G$ and $w=(P/P_i^k)(\partial)u$. If $k\le e_i-f_i$, then
$G\mid P/P_i^k$, say $P/P_i^k=GH$; for $h$ in a power $\mathfrak n^{\ell}\subseteq J$ of the ideal of $\Var(J)$,
$h\cdot P/P_i^k=G(Hh)\in I$, so $\Ann(w)\supseteq\mathfrak n^{\ell}$ and $w\in\EP$. If $k>e_i-f_i$ and $w\in\EP$, pick
$h\in\Ann(w)$ with $P_i\nmid h$ (as in the proof of \cref{thm:ident}); then $hP/P_i^k\in I$, so $G\mid hP/P_i^k$, but
$P_i$ occurs in $hP/P_i^k$ with multiplicity $e_i-k<f_i$, a contradiction. Hence
$f_i=e_i-\max\{k:\ w\in\EP\}$.
\end{proof}

\subsection{Decomposition and stability}
\begin{proof}[Proof of \cref{thm:factor}]
We prove the statement for open connected $\Omega$ with $H^1_{dR}(\Omega)=0$, which holds for simply connected domains.

\emph{Step 1: finite modules have no obstruction.} Let $M$ be an $\ring$-module of finite dimension over $\C$. Then
$\operatorname{Ext}^1_\ring(M,\mathcal A)=0$ and $\dim\operatorname{Hom}_\ring(M,\mathcal A)=\dim M$. By the long exact sequence it is
enough to check one-dimensional modules $\ring/\mathfrak m_z$, since $M$ has a composition series with such
quotients. A free resolution of $\ring/\mathfrak m_z$ is the Koszul complex of $(\tau-\tau_z,\xi-\xi_z)$, so
$\operatorname{Ext}^\bullet(\ring/\mathfrak m_z,\mathcal A)$ is the cohomology of
$\mathcal A\to\mathcal A^2\to\mathcal A$, $v\mapsto\big((\partial_t-\tau_z)v,(\partial_x-\xi_z)v\big)$, $(a,b)\mapsto(\partial_t-\tau_z)b-(\partial_x-\xi_z)a$.
Multiplication by $\psi_{-z}$ turns it into the de Rham complex of $\Omega$ with coefficients in $\mathcal A$, whose first
cohomology is $H^1_{dR}(\Omega)=0$ (for distributions by de Rham's theorem), and whose zeroth cohomology is the constants,
of dimension one.

\emph{Step 2: two factors.} For coprime $Q_1,Q_2$, $(Q_1)\cap(Q_2)=(Q_1Q_2)$ by unique factorization, and
$0\to\ring/(Q_1Q_2)\to\ring/(Q_1)\oplus\ring/(Q_2)\to\ring/(Q_1,Q_2)\to0$ is exact ($1\mapsto(1,1)$, $(a,b)\mapsto a-b$).
Applying $\operatorname{Hom}_\ring(\cdot,\mathcal A)$, and using $\operatorname{Hom}(\ring/(f),\mathcal A)=\ker_{\mathcal A}f$,
gives the exact sequence
$0\to\operatorname{Hom}(\ring/(Q_1,Q_2),\mathcal A)\to\ker Q_1\oplus\ker Q_2\xrightarrow{(u_1,u_2)\mapsto u_1+u_2}\ker Q_1Q_2\to
\operatorname{Ext}^1(\ring/(Q_1,Q_2),\mathcal A)$. The quotient $\ring/(Q_1,Q_2)$ is finite-dimensional by
\cref{lem:gcd}, so by Step~1 the last group vanishes: the sum map is onto, and its kernel has dimension
$\dim\ring/(Q_1,Q_2)$ and consists of solutions of a zero-dimensional ideal, i.e.\ of exponential-polynomials.

\emph{Step 3: induction.} $P'=Q_2\cdots Q_N$ is coprime to $Q_1$, so $\ker P=\ker Q_1+\ker P'$, and by induction
$\ker P'=\sum_{j\ge2}\ker Q_j$. Dimensions of kernels add along the composition of the two sum maps, and
$\dim\ring/(Q_1,P')=\sum_{j\ge2}\dim\ring/(Q_1,Q_j)$ because $\ring/(Q_1,Q_2Q_3\cdots)$ is the direct sum of its
localizations at the finitely many common zeros, where the intersection multiplicities add. This gives
\eqref{eq:deltaP}.
\end{proof}

\begin{proof}[Proof of \cref{thm:hm}]
The sum operator $\sumop$ is bounded from $\bigoplus_jM_j$ into $\ker_{\Hs}P$, and every $u$ in its range has
components in $\Hs$. Its range is all of $\ker_{\Hs}P$ iff the range is closed and dense. By the closed range theorem the
range is closed iff $\sumop$ is bounded below on $(\ker\sumop)^\perp$, i.e.\ iff $\gN>0$, and then the minimum-norm
preimage satisfies \eqref{eq:stability}. For two closed subspaces, $M_1+M_2$ is closed iff $\cF<1$, and
$\gamma_2=1-\cF$~\cite{Deu01}. In finite dimension every subspace sum is closed, and the sum equals the span of
the sampled exponentials of $P$, each of which lies in the kernel of one factor.
\end{proof}

\begin{proof}[Proof of \cref{prop:discont}]
For real $\zeta>c/\nu$ let $u_\zeta=e^{\lambda_\zeta t+\zeta x}$, $\lambda_\zeta=\nu\zeta^2-c\zeta>0$; then $Q(\partial)u_\zeta=0$ and
$\norm{u_\zeta}_{\Hs(\Omega)}\le C_s(1+\zeta)^{2s}e^{\lambda_\zeta T}\norm{e^{\zeta x}}_{L^2(a,b)}$. With base time $t_0$ and horizon
$t$, \cref{thm:eint} gives
\[
  \Eint=e^{\lambda_\zeta t_0}\big(e^{-c\zeta t}-1\big)\big(e^{\nu\zeta^2t}-1\big)\,e^{\zeta x},
\]
whose $L^2(\Omega)$ norm is at least a constant times $e^{\nu\zeta^2t}\norm{e^{\zeta x}}_{L^2(a,b)}$ for large $\zeta$. With
$t_0=0$ and $t=T$ the ratio therefore grows at least like
\[
  e^{\nu\zeta^2T-\lambda_\zeta T}(1+\zeta)^{-2s}=e^{c\zeta T}(1+\zeta)^{-2s},
\]
which is unbounded for every $s$. For the second statement, $v+w$ is killed by $Q$ and by
$R=(\tau+c\xi)(\tau-\nu\xi^2)$. At a common zero of $Q$ and $R$, either $\tau=-c\xi$, and then $Q=0$ gives
$\nu\xi^2=0$, or $\tau=\nu\xi^2$, and then $c\xi=0$; in both cases $(\tau,\xi)=(0,0)$. So $(Q,R)$ has a single zero, and
its common kernel is a finite-dimensional space of exponential-polynomials (\cref{lem:ep}).
\end{proof}
